\chapter{Lijnintegralen en meervoudige integralen}\label{ch:b2:multint}

Dit hoofdstuk breidt de integratie uit van intervallen naar krommen
en naar gebieden van het vlak en van de ruimte. Lijnintegralen
integreren een \emph{\omterm{def:b2:multint:lineint}{differentiaalvorm}} $P\,\dd x + Q\,\dd y$ langs
een georiënteerde boog; dubbele en drievoudige integralen
integreren functies over twee- en driedimensionale gebieden. De
twee theorieën ontmoeten elkaar in de \emph{stelling van
Green--Riemann}, de tweedimensionale hoofdstelling van de
integraalrekening, en het voornaamste rekengereedschap is overal de
\emph{formule voor de verandering van veranderlijken}, waarvan de
vervormingsfactor de absolute waarde van de jacobideterminant is.

\section{Lijnintegralen}

\begin{definition}[Differentiaalvorm; lijnintegraal]\label{def:b2:multint:lineint}
Zij $U \subseteq \R^2$ \omterm{def:b2:metric:topology}{open}. Een
\emph{differentiaalvorm}\index{differentiaalvorm} van graad $1$ en
van klasse $\mathcal{C}^0$ op $U$ is een uitdrukking $\omega =
P\,\dd x + Q\,\dd y$ met $P, Q \colon U \to \R$ \omterm{def:b2:metric:continuity}{continu} ---
formeel een \omterm{def:b2:metric:continuity}{continue} afbeelding van $U$ in de \omterm{def:b2:linalg:dual}{duale ruimte} van
$\R^2$, $\omega(M) = P(M)\,e_1^* + Q(M)\,e_2^*$. Voor een
$\mathcal{C}^1$-boog $\gamma \colon [a, b] \to U$, $\gamma(t) =
(x(t), y(t))$, is de \emph{lijnintegraal}\index{lijnintegraal} van
$\omega$ langs $\gamma$
\[
\int_\gamma \omega
= \int_a^b \Bigl(P(\gamma(t))\,x'(t)
+ Q(\gamma(t))\,y'(t)\Bigr)\,\dd t .
\]
De definities breiden woordelijk uit tot $\R^3$ (vormen $P\,\dd x +
Q\,\dd y + R\,\dd z$) en tot stuksgewijs $\mathcal{C}^1$-bogen (som
over de stukken).
\end{definition}

\begin{proposition}[Invariantie en oriëntatie]\label{prop:b2:multint:lineinv}
De \omterm{def:b2:multint:lineint}{lijnintegraal} verandert niet onder een stijgende
$\mathcal{C}^1$-parameterverandering, en verandert van teken onder
een dalende. Zij hangt dus alleen van de \emph{georiënteerde}
\omterm{def:b2:curves:reparam}{meetkundige boog} af.
\end{proposition}

\begin{proof}
Is $\theta \colon [c, d] \to [a, b]$ een \omterm{def:b2:curves:reparam}{parameterverandering} en
$\tilde\gamma = \gamma \circ \theta$, dan geldt volgens de
kettingregel en de verandering van veranderlijke $t = \theta(u)$ in
één veranderlijke
\[
\int_{\tilde\gamma}\omega
= \int_c^d \bigl(P(\gamma(\theta(u)))\,x'(\theta(u))
+ Q(\gamma(\theta(u)))\,y'(\theta(u))\bigr)\,\theta'(u)\,\dd u
= \pm\int_a^b \bigl(Px' + Qy'\bigr)(t)\,\dd t ,
\]
met teken $+$ als $\theta$ stijgt ($\theta(c) = a$) en $-$ als zij
daalt (de grenzen verwisselen).
\end{proof}

\begin{example}[Arbeid van een kracht; circulatie]\label{ex:b2:multint:work}
Is $F = (P, Q)$ een krachtveld, dan is $\int_\gamma P\dd x + Q\dd y
= \int_a^b \langle F(\gamma(t)), \gamma'(t)\rangle\,\dd t$ de
\emph{arbeid} van $F$ langs $\gamma$. Voor $\omega = -y\,\dd x +
x\,\dd y$ langs de eenheidscirkel tegen de wijzers van de klok in,
$\gamma(t) = (\cos t, \sin t)$:
\[
\int_\gamma \omega
= \int_0^{2\pi}\bigl((-\sin t)(-\sin t)
+ \cos t\cos t\bigr)\,\dd t = 2\pi ,
\]
tweemaal de omsloten \omterm{def:b2:multint:domain}{oppervlakte} --- een eerste hint van
Green--Riemann.
\end{example}

\begin{example}[Eén integraal, twee parametriseringen, één
tekenval]\label{ex:b2:multint:semicircle}
Bereken $\int_\gamma x\,\dd y$ langs de bovenste
eenheidshalfcirkel van $(1, 0)$ naar $(-1, 0)$. Met $\gamma(t) =
(\cos t, \sin t)$, $t \in \intcc0\pi$:
\[
\int_0^\pi\cos t\cdot\cos t\,\dd t = \frac\pi2 .
\]
Met de grafiekparametrisering $x \mapsto (x, \sqrt{1 - x^2})$, $x$
van $1$ naar $-1$ (let op de richting!):
\[
\int_1^{-1}x\cdot\frac{-x}{\sqrt{1 - x^2}}\,\dd x
= \int_{-1}^{1}\frac{x^2}{\sqrt{1 - x^2}}\,\dd x
= \frac\pi2
\]
($x = \sin u$ herleidt haar tot een integraal van Wallis). Dezelfde
waarde, zoals \cref{prop:b2:multint:lineinv} waarborgt --- maar
alleen omdat beide doorlopen van $(1,0)$ naar $(-1,0)$ gaan; de
reis omkeren keert het teken om. Het pad sluiten langs de $x$-as
(waar $\dd y = 0$) voegt niets toe, en het totaal $\frac\pi2$ is de
\omterm{def:b2:multint:domain}{oppervlakte} van de halve schijf: het eerste geval van de
randformules voor \omterm{def:b2:surfaces:area}{oppervlakten} van Green--Riemann hieronder.
\end{example}

\begin{definition}[Exacte en gesloten vormen]\label{def:b2:multint:exact}
De vorm $\omega = P\,\dd x + Q\,\dd y$ van klasse $\mathcal{C}^0$
heet \emph{exact}\index{exacte vorm} op $U$ wanneer er een $f \in
\mathcal{C}^1(U)$ bestaat (een
\emph{potentiaal}\index{potentiaal}) met $\omega = \dd f$, dat wil
zeggen $P = f_x$ en $Q = f_y$. Een $\mathcal{C}^1$-vorm heet
\emph{gesloten}\index{gesloten vorm} wanneer $P_y = Q_x$ op $U$.
\end{definition}

\begin{theorem}[Hoofdstelling voor
lijnintegralen]\label{thm:b2:multint:ftc}
Is $\omega = \dd f$ \omterm{def:b2:multint:exact}{exact} en $\gamma$ een stuksgewijs
$\mathcal{C}^1$-boog in $U$ van $A$ naar $B$, dan is
\[
\int_\gamma \omega = f(B) - f(A) .
\]
In het bijzonder is de integraal van een \omterm{def:b2:multint:exact}{exacte vorm} langs elke
\omterm{def:b2:multint:exact}{gesloten} boog nul, en is elke exacte $\mathcal{C}^1$-vorm \omterm{def:b2:multint:exact}{gesloten}.
\end{theorem}

\begin{proof}
$\frac{\dd}{\dd t}f(\gamma(t)) = f_x(\gamma(t))x'(t) +
f_y(\gamma(t))y'(t)$ volgens de kettingregel
(\cref{ch:b2:diffcalc}), dus is de integrand in
\cref{def:b2:multint:lineint} de afgeleide van $t \mapsto
f(\gamma(t))$, en de hoofdstelling van de integraalrekening geeft
het resultaat op elk stuk; de tussenwaarden telescoperen. Dat
exacte $\mathcal{C}^1$-vormen \omterm{def:b2:multint:exact}{gesloten} zijn, is de stelling van
Schwarz: $P_y = f_{xy} = f_{yx} = Q_x$.
\end{proof}

\begin{example}[Een potentiaal reconstrueren]\label{ex:b2:multint:potential}
Zij $\omega = y\,\eu^{xy}\,\dd x + (x\,\eu^{xy} + 2y)\,\dd y$ op
$\R^2$. Zij is \omterm{def:b2:multint:exact}{gesloten}: beide gemengde afgeleiden zijn gelijk aan
$\eu^{xy}(1 + xy)$. Om een \omterm{def:b2:multint:exact}{potentiaal} te vinden, integreer je $P$
naar $x$ bij vaste $y$:
\[
f(x, y) = \int y\,\eu^{xy}\,\dd x = \eu^{xy} + c(y),
\]
en pas je $c$ aan door $f_y$ te laten kloppen: $x\,\eu^{xy} + c'(y)
= x\,\eu^{xy} + 2y$ geeft $c(y) = y^2$. Dus $f(x,y) = \eu^{xy} +
y^2$, en voor elke stuksgewijs $\mathcal C^1$-boog van $(0,0)$ naar
$(1,1)$ is
\[
\int_\gamma\omega = f(1,1) - f(0,0) = (\eu + 1) - 1 = \eu ,
\]
onafhankelijk van het pad --- het recept in twee stappen
(integreren naar $x$, corrigeren in $y$) is de praktische omkering
van \cref{thm:b2:multint:ftc} op gebieden waar gesloten vormen
\omterm{def:b2:multint:exact}{exact} zijn.
\end{example}

\begin{example}[Gesloten impliceert niet exact]\label{ex:b2:multint:angleform}
Op $U = \R^2 \setminus \{0\}$ is de \emph{hoekvorm}
\[
\omega = \frac{-y\,\dd x + x\,\dd y}{x^2 + y^2}
\]
\omterm{def:b2:multint:exact}{gesloten} (rechtstreekse berekening: zowel $P_y$ als $Q_x$ is gelijk
aan $\frac{y^2 - x^2}{(x^2+y^2)^2}$), maar haar integraal langs de
eenheidscirkel is $2\pi \neq 0$ (dezelfde berekening als in
\cref{ex:b2:multint:work}, gedeeld door $1$): $\omega$ is niet
\omterm{def:b2:multint:exact}{exact} op $U$. Lokaal is $\omega = \dd\theta$ voor een bepaling
$\theta$ van de poolhoek; het falen is globaal --- de hoek kan niet
\omterm{def:b2:metric:continuity}{continu} rond het gat worden gedefinieerd. Op gebieden zonder gaten
verdwijnt de pathologie: op een \emph{stervormige} \omterm{def:b2:metric:topology}{open} verzameling
is elke \omterm{def:b2:multint:exact}{gesloten} $\mathcal{C}^1$-vorm \omterm{def:b2:multint:exact}{exact} (het lemma van
Poincaré, \cref{exo:b2:multint:8}).
\end{example}

\section{Dubbele integralen}

Wij nemen de theorie van de riemann-integraal in één veranderlijke
als bekend aan (volume van bachelorjaar 1, en
\cref{ch:b2:integration}) en schetsen haar versie met twee
veranderlijken. Een functie $f$ die \omterm{def:b2:metric:continuity}{continu} is op een rechthoek $R
= [a, b] \times [c, d]$ heeft een dubbele integraal $\iint_R f$,
gedefinieerd met Riemannsommen over roosters precies als in één
veranderlijke, en berekend door iteratie:

\begin{theorem}[Fubini op een rechthoek]\label{thm:b2:multint:fubini}
Voor $f$ \omterm{def:b2:metric:continuity}{continu} op $R = [a,b] \times [c,d]$ geldt
\[
\iint_R f
= \int_a^b \Bigl(\int_c^d f(x, y)\,\dd y\Bigr)\dd x
= \int_c^d \Bigl(\int_a^b f(x, y)\,\dd x\Bigr)\dd y .
\]
\end{theorem}

\begin{proof}
Zet $F(x) = \int_c^d f(x, y)\,\dd y$. De uniforme \omterm{def:b2:metric:continuity}{continuïteit} van
$f$ op het \omterm{def:b2:metric:compact}{compacte} $R$ maakt $F$ \omterm{def:b2:metric:continuity}{continu} (gedomineerde schatting:
$\abs{F(x) - F(x')} \leq (d - c)\sup_y\abs{f(x,y) - f(x',y)}$).
Verdeel nu $[a,b]$ en $[c,d]$ elk in $n$ gelijke delen, wat een
rooster van cellen $R_{ij}$ met \omterm{def:b2:multint:domain}{oppervlakte} $\Delta x\,\Delta y$
geeft. Op elke cel is $\inf_{R_{ij}} f \cdot \Delta x \Delta y \leq
\int_{x_{i-1}}^{x_i}\int_{y_{j-1}}^{y_j} f(x,y)\,\dd y\,\dd x \leq
\sup_{R_{ij}} f \cdot \Delta x \Delta y$ wegens de monotonie van de
integraal in één veranderlijke (tweemaal toegepast). Sommeren over
de cellen klemt de herhaalde integraal $\int_a^b F$ tussen de
onder- en bovensommen van Riemann van het rooster; wegens de
uniforme \omterm{def:b2:metric:continuity}{continuïteit} convergeren beide sommen naar de
gemeenschappelijke waarde die $\iint_R f$ definieert als $n \to
\infty$. Hetzelfde argument werkt met de rollen van $x$ en $y$
verwisseld, dus beide herhaalde integralen zijn gelijk aan $\iint_R
f$.
\end{proof}

\begin{remark}
De \omterm{def:b2:metric:continuity}{continuïteit} op een \omterm{def:b2:metric:compact}{compacte} rechthoek doet in het bewijs van
Fubini echt werk: zij levert de uniforme \omterm{def:b2:metric:continuity}{continuïteit} die de
Riemannsommen inklemt. Voor wildere integranden faalt de uitspraak
werkelijk --- er bestaan functies waarvan de twee herhaalde
integralen bestaan en verschillen. De eerlijke algemene stelling,
met integreerbaarheid als enige hypothese, is de stelling van
Fubini voor de lebesgue-integraal, bewezen in het volume van
bachelorjaar 3; alles in dit hoofdstuk blijft binnen het \omterm{def:b2:metric:continuity}{continue}
kader waar het elementaire bewijs hierboven \omterm{def:b2:metric:complete}{volledig} is.
\end{remark}

\begin{definition}[Elementaire gebieden]\label{def:b2:multint:domain}
Een gebied $D \subseteq \R^2$ heet \emph{$y$-elementair} wanneer
\[
D = \{(x, y) : a \leq x \leq b,\
\varphi_1(x) \leq y \leq \varphi_2(x)\}
\]
met $\varphi_1 \leq \varphi_2$ \omterm{def:b2:metric:continuity}{continu} op $[a,b]$
($x$-elementair: \omterm{def:b2:quadratic:adjoint}{symmetrisch}). Voor $f$ \omterm{def:b2:metric:continuity}{continu} op een
$y$-elementair gebied $D$ is
\[
\iint_D f
= \int_a^b\Bigl(
\int_{\varphi_1(x)}^{\varphi_2(x)} f(x,y)\,\dd y\Bigr)\dd x ,
\]
en men gaat na (door $f$ met een benaderingsargument uit te
breiden, of door onder te verdelen) dat wanneer $D$ in beide
richtingen elementair is, de twee herhaalde integralen
overeenstemmen. Gebieden die in eindig veel elementaire stukken
worden gesneden, worden met de additiviteit behandeld, en de
\emph{oppervlakte}\index{oppervlakte!van een vlak gebied} van $D$
is $\operatorname{Area}(D) = \iint_D 1$.
\end{definition}

\begin{example}\label{ex:b2:multint:triangle}
Op de driehoek $D = \{0 \leq x \leq 1,\ 0 \leq y \leq x\}$:
\[
\iint_D xy \,\dd x\,\dd y
= \int_0^1 x\Bigl(\int_0^x y\,\dd y\Bigr)\dd x
= \int_0^1 x\cdot\frac{x^2}{2}\,\dd x = \frac18 .
\]
De volgorde verwisselen ($x$ van $y$ tot $1$): $\int_0^1
y\bigl(\int_y^1 x\,\dd x\bigr)\dd y = \int_0^1 y\,\frac{1 -
y^2}{2}\,\dd y = \frac18$ --- dezelfde waarde, een andere
berekening: de volgorde van integreren goed kiezen is de helft van
het vak.
\end{example}

\begin{example}[Wanneer maar één volgorde werkt]\label{ex:b2:multint:swaporder}
Bereken $I = \displaystyle\int_0^1\!\!\int_x^1
\eu^{y^2}\,\dd y\,\dd x$. Zoals zij er staat, heeft de binnenste
integraal $\int\eu^{y^2}\dd y$ geen elementaire primitieve: de
berekening loopt vast. Maar het gebied is de driehoek $0 \leq x
\leq y \leq 1$, die in beide richtingen elementair is; de volgorde
verwisselen geeft
\[
I = \int_0^1\!\!\int_0^y \eu^{y^2}\,\dd x\,\dd y
= \int_0^1 y\,\eu^{y^2}\,\dd y
= \Bigl[\tfrac12\eu^{y^2}\Bigr]_0^1 = \frac{\eu - 1}{2} .
\]
De binnenste veranderlijke $x$ kwam in de integrand nergens voor,
dus haar eerst integreren bracht precies de factor $y$ voort die de
buitenste integraal onmiddellijk maakt. Moraal: Fubini is niet
alleen een vergunning om te itereren --- het is een vergunning om
te \emph{kiezen}, en de juiste volgorde kan van een onmogelijke
integraal een regel van één zin maken. Schets altijd het gebied en
lees beide beschrijvingen af voordat je begint.
\end{example}

\begin{theorem}[Verandering van veranderlijken]\label{thm:b2:multint:changevar}
Zij $\Phi \colon U' \to U$ een $\mathcal{C}^1$-diffeomorfisme
tussen \omterm{def:b2:metric:topology}{open} verzamelingen van $\R^2$, zij $K \subseteq U$ een
\omterm{def:b2:metric:compact}{compact} gebied dat in elementaire stukken is gesneden met $K' =
\Phi^{-1}(K)$, en zij $f$ \omterm{def:b2:metric:continuity}{continu} op $K$. Dan is
\[
\iint_K f(x, y)\,\dd x\,\dd y
= \iint_{K'} f\bigl(\Phi(u, v)\bigr)\,
\abs{\det J_\Phi(u, v)}\,\dd u\,\dd v .
\]
\end{theorem}
\admitted

\begin{remark}
Het volledige bewijs --- $\Phi$ door haar \omterm{def:b2:diffcalc:differential}{differentiaal} benaderen
op een fijn rooster en de randcellen beheersen --- is lang maar
niet diep; het wordt \omterm{def:b2:metric:complete}{volledig} uitgevoerd in de maattheorie van
bachelorjaar 3, als gevolg van de theorie van Lebesgue. De
heuristiek is het beeld dat al voor de \omterm{def:b2:multint:domain}{oppervlakte} van een oppervlak
werd gebruikt: een klein vierkant met zijde $\dd u$ in $(u, v)$
wordt tot op eerste orde afgebeeld op het parallellogram
opgespannen door $\Phi_u\,\dd u$ en $\Phi_v\,\dd v$, waarvan de
\omterm{def:b2:multint:domain}{oppervlakte} $\abs{\det J_\Phi}\,\dd u\,\dd v$ is
(\cref{lem:b2:surfaces:lagrange}).
\end{remark}

\begin{remark}[Methode: de verandering van veranderlijken kiezen]
Drie reflexen dekken de meeste gevallen. \emph{Symmetrie van de
integrand}: $x^2 + y^2$ roept om poolcoördinaten, een
productstructuur roept om het behouden van de cartesische assen.
\emph{Vorm van de rand}: randen $u(x,y) = c_1$, $v(x,y) = c_2$
smeken om de coördinaten $(u, v)$ zelf, zoals in het voorbeeld met
het hyperbolische gebied hieronder --- het gebied wordt een
rechthoek, en dat is de hele overwinning. \emph{Lineaire
structuur}: uitdrukkingen in $x + y$ en $x - y$ nodigen uit tot de
rotatie over $45$ graden of tot een afschuiving
(\cref{ex:b2:multint:affine}). In alle gevallen zijn er drie hokjes
aan te vinken voordat je integreert: de afbeelding is een bijectie
van het nieuwe gebied op het oude; haar jacobideterminant wordt
berekend \emph{in de richting die werkelijk wordt gebruikt}
(inverteer op het einde als dat gemakkelijker is); en de
jacobideterminant komt binnen met haar absolute waarde.
\end{remark}

\begin{example}[Poolcoördinaten]\label{ex:b2:multint:polar}
$\Phi(\rho, \alpha) = (\rho\cos\alpha,\ \rho\sin\alpha)$ heeft
\[
J_\Phi = \begin{pmatrix}
\cos\alpha & -\rho\sin\alpha\\
\sin\alpha & \rho\cos\alpha
\end{pmatrix},
\qquad \det J_\Phi = \rho ,
\]
dus $\dd x\,\dd y = \rho\,\dd\rho\,\dd\alpha$. Voor de schijf $D_R$
met straal $R$:
\[
\iint_{D_R} e^{-(x^2 + y^2)}\,\dd x\,\dd y
= \int_0^{2\pi}\!\!\int_0^R e^{-\rho^2}\rho\,\dd\rho\,\dd\alpha
= \pi\bigl(1 - e^{-R^2}\bigr)
\xrightarrow[R\to\infty]{} \pi .
\]
Vergelijken met het vierkant $[-R, R]^2$ (dat tussen de schijven
$D_R$ en $D_{R\sqrt2}$ geklemd zit, met overal positieve
integranden) geeft $\bigl(\int_{-\infty}^\infty
e^{-x^2}\dd x\bigr)^2 = \pi$:
\[
\boxed{\ \int_{-\infty}^{+\infty} e^{-x^2}\,\dd x = \sqrt{\pi}\ }
\]
--- opnieuw de integraal van Gauss, nu met haar beroemdste bewijs
(vergelijk de afleiding in één veranderlijke in
\cref{ch:b2:integration}).
\end{example}

\begin{example}[Affiene veranderingen van veranderlijken]\label{ex:b2:multint:affine}
Voor een \omterm{def:b2:affine:subspace}{affiene afbeelding} $\Phi(u, v) = M(u, v)^{\mathsf T} + C$
met $M$ inverteerbaar is de \omterm{def:b2:diffcalc:differential}{jacobimatrix} de constante matrix $M$:
\omterm{def:b2:surfaces:area}{oppervlakten} worden vermenigvuldigd met de constante factor
$\abs{\det M}$ --- de belofte uit \cref{ch:b2:affine} is nu een
stelling. Twee onmiddellijke toepassingen. De ellips
$\frac{x^2}{a^2} + \frac{y^2}{b^2} \leq 1$ is het beeld van de
eenheidsschijf onder $(u, v) \mapsto (au, bv)$, dus is haar
\omterm{def:b2:multint:domain}{oppervlakte} $ab \cdot \pi$ --- zonder berekening. En voor de
integraal van $f(x + y)$ over het vierkant $K = \intcc01^2$ maakt
de afschuiving $\Phi(u, v) = (u - v, v)$ (\omterm{def:b2:linalg:det}{determinant} $1$) er een
integraal van $f(u)$ over een parallellogram van, dat Fubini bij
constante $u$ in schijven snijdt: met $f = \exp$,
\[
\iint_K \eu^{x+y}\,\dd x\,\dd y
= \Bigl(\int_0^1 \eu^x\,\dd x\Bigr)^2 = (\eu - 1)^2,
\]
zoals de productstructuur bevestigt. Coördinaten kiezen die bij de
integrand passen --- niet bij het gebied --- is de andere helft van
het vak.
\end{example}


\begin{example}[Coördinaten aangepast aan een kromlijnig
gebied]\label{ex:b2:multint:hyperbolicregion}
Zij $D$ het gebied in het eerste kwadrant begrensd door de
hyperbolen $xy = 1$ en $xy = 3$ en de rechten $y = x$ en $y = 3x$.
In de coördinaten $u = xy$, $v = y/x$ wordt het gebied het vierkant
$\intcc13 \times \intcc13$; inverteren geeft
\[
x = \sqrt{u/v}, \qquad y = \sqrt{uv},
\qquad
\det J = x_uy_v - x_vy_u = \frac{1}{2v}
\]
(een berekening van twee regels met $x = u^{1/2}v^{-1/2}$, $y =
u^{1/2}v^{1/2}$). Bijgevolg is
\[
\operatorname{Area}(D)
= \int_1^3\!\!\int_1^3\frac{\dd u\,\dd v}{2v}
= 2\cdot\frac{\ln 3}{2} = \ln 3 \approx 1.10 .
\]
$D$ in cartesische coördinaten in schijven snijden betekent haar in
drie stukken met hyperbolische en lineaire randen snijden ---
haalbaar, vreugdeloos en foutgevoelig. De moraal herhaalt
\cref{ex:b2:multint:affine} in volle kracht: lees de
randvergelijkingen, en laat \emph{hen} de coördinaten kiezen; de
jacobideterminant zet dan de celoppervlakte van het kromlijnige
rooster om, precies zoals $\rho$ dat voor poolcoördinaten deed.
\end{example}

\begin{example}[Gemiddelde waarden]\label{ex:b2:multint:average}
De \emph{gemiddelde waarde} van $f$ over een gebied $D$ is
$\frac1{\operatorname{Area}(D)}\iint_Df$. Voorbeeld: de gemiddelde
afstand tot het middelpunt voor een uniform gekozen punt van de
schijf met straal $R$ is
\[
\frac{1}{\pi R^2}\int_0^{2\pi}\!\!\int_0^R
\rho\cdot\rho\,\dd\rho\,\dd\alpha
= \frac{2\pi R^3/3}{\pi R^2} = \frac{2R}3 ,
\]
niet $R/2$: een uniforme \omterm{def:b2:multint:domain}{oppervlakte} legt meer massa bij grote
stralen (de ring op straal $\rho$ heeft gewicht evenredig met
$\rho$), zodat het gemiddelde voorbij het halverwegepunt ligt. Deze
factor juist krijgen is precies de jacobiaan van de poolcoördinaten
aan het werk, en dezelfde weging verklaart het zwaartepunt $\bar z
= 3R/8$ van de halve bol dat verderop in het hoofdstuk wordt
berekend, in plaats van $R/2$.
\end{example}

\section{De stelling van Green--Riemann}

\begin{theorem}[Green--Riemann]\label{thm:b2:multint:green}
Zij $K \subseteq \R^2$ een \omterm{def:b2:metric:compact}{compact} gebied dat in beide richtingen
elementair is (of een eindige vereniging van zulke gebieden, langs
segmenten aaneengelijmd), met als rand $\partial K$ een
stuksgewijs $\mathcal{C}^1$ \omterm{def:b2:multint:exact}{gesloten} kromme, georiënteerd
\emph{tegen de wijzers van de klok in} (het gebied blijft links).
Voor $P, Q$ van klasse $\mathcal{C}^1$ op een omgeving van $K$
geldt
\[
\oint_{\partial K} P\,\dd x + Q\,\dd y
= \iint_K \Bigl(\frac{\partial Q}{\partial x}
- \frac{\partial P}{\partial y}\Bigr)\,\dd x\,\dd y .
\]
\end{theorem}

\begin{proof}
Ten eerste zijn beide leden additief wanneer $K$ langs een segment
in twee stukken $K_1, K_2$ wordt gesneden: de dubbele integralen
tellen op wegens de additiviteit van $\iint$; en wat de
randintegralen betreft, doorlopen de randen van $K_1$ en $K_2$
tegen de wijzers van de klok in de inwendige snede elk eenmaal, in
\emph{tegengestelde} richtingen, zodat in de som
\[
\oint_{\partial K_1} + \oint_{\partial K_2}
= \oint_{\partial K} + (\text{de snede, in beide richtingen})
= \oint_{\partial K},
\]
de twee passages langs de snede elkaar opheffen
(\cref{prop:b2:multint:lineinv}) en alleen de buitenrand
overblijft. Door eindig veel sneden te herhalen, volstaat het een
\omterm{def:b2:multint:domain}{elementair gebied} te behandelen. Wij bewijzen $\oint P\,\dd x =
-\iint_K P_y$ op een $y$-elementair gebied $D = \{a \leq x \leq
b,\ \varphi_1(x) \leq y \leq \varphi_2(x)\}$; de identiteit $\oint
Q\,\dd y = \iint_K Q_x$ is \omterm{def:b2:quadratic:adjoint}{symmetrisch} ($x$-elementair), en de
stelling is hun som.

Bereken de dubbele integraal met Fubini en de hoofdstelling in één
veranderlijke:
\[
\iint_D \frac{\partial P}{\partial y}\,\dd x\,\dd y
= \int_a^b \bigl(P(x, \varphi_2(x)) - P(x, \varphi_1(x))\bigr)
\,\dd x .
\]
Nu bestaat de rand van $D$, tegen de wijzers van de klok in, uit:
de onderste grafiek $y = \varphi_1(x)$ van links naar rechts
doorlopen, het rechtse verticale segment $x = b$ (omhoog), de
bovenste grafiek $y = \varphi_2(x)$ \emph{van rechts naar links}
doorlopen, en het linkse verticale segment $x = a$ (omlaag). Langs
de verticale segmenten is $x$ constant, dus dragen zij $0$ bij aan
$\oint P\,\dd x$; de grafieken, geparametriseerd met $x$, geven
\[
\oint_{\partial D} P\,\dd x
= \int_a^b P(x, \varphi_1(x))\,\dd x
- \int_a^b P(x, \varphi_2(x))\,\dd x
= -\iint_D \frac{\partial P}{\partial y}\,\dd x\,\dd y .
\qedhere
\]
\end{proof}

\begin{corollary}[Oppervlakte via de rand]\label{cor:b2:multint:area}
Onder de hypothesen van \cref{thm:b2:multint:green} is
\[
\operatorname{Area}(K)
= \oint_{\partial K} x\,\dd y
= -\oint_{\partial K} y\,\dd x
= \frac12\oint_{\partial K} x\,\dd y - y\,\dd x .
\]
\end{corollary}

\begin{proof}
Pas Green--Riemann toe op $(P, Q) = (0, x)$, $(-y, 0)$ en
$\frac12(-y, x)$: telkens is $Q_x - P_y = 1$.
\end{proof}

\begin{remark}[Kiezen tussen de drie oppervlakteformules]
De drie randformules zijn gelijk, maar in de praktijk niet
onderling verwisselbaar. Gebruik $\oint x\,\dd y$ wanneer de
parametrisering $\dd y$ eenvoudig maakt (grafieken boven de
$y$-as), $-\oint y\,\dd x$ \omterm{def:b2:quadratic:adjoint}{symmetrisch}, en de \omterm{def:b2:quadratic:adjoint}{symmetrische} halve
som wanneer de parametrisering $x$ en $y$ gelijk behandelt --- voor
de ellips leverde zij een constante integrand op, zonder enige
goniometrische linearisatie. Op veelhoekige randen wordt de halve
som de veterformule van \cref{exo:b2:multint:12}, het algoritme van
de landmeters. En wordt de rand door de gegeven parametrisering met
de wijzers van de klok mee doorlopen, dan geven alle drie de
formules \emph{min} de \omterm{def:b2:multint:domain}{oppervlakte}: een negatief resultaat is geen
rekenfout maar een melding over de oriëntatie --- keer het teken
om, of de parametrisering.
\end{remark}

\begin{example}[Oppervlakte van de ellips]\label{ex:b2:multint:ellipsearea}
Voor $x = a\cos t$, $y = b\sin t$, $t \in [0, 2\pi]$:
\[
\operatorname{Area}
= \frac12\int_0^{2\pi}\bigl(a\cos t \cdot b\cos t
- b\sin t\cdot(-a\sin t)\bigr)\,\dd t
= \frac{ab}{2}\int_0^{2\pi}\dd t = \pi ab .
\]
\end{example}

\begin{example}[Green--Riemann als
kruiscontrole]\label{ex:b2:multint:greencheck}
Neem $P = -y^3$, $Q = x^3$ op de \omterm{def:b2:multint:exact}{gesloten} eenheidsschijf $D$. De
randkant, met $\gamma(t) = (\cos t, \sin t)$:
\[
\oint_{\partial D}P\,\dd x + Q\,\dd y
= \int_0^{2\pi}\bigl(\sin^4 t + \cos^4 t\bigr)\dd t
= 2\pi\cdot\Bigl(\frac38 + \frac38\Bigr) = \frac{3\pi}2 ,
\]
door lineariseren ($\sin^4 + \cos^4 = \tfrac34 +
\tfrac14\cos4t$). De inwendige kant:
\[
\iint_D(Q_x - P_y)\,\dd x\,\dd y
= \iint_D 3(x^2 + y^2)\,\dd x\,\dd y
= 3\int_0^{2\pi}\!\!\int_0^1\rho^3\,\dd\rho\,\dd\alpha
= \frac{3\pi}2 .
\]
Hetzelfde getal, twee heel verschillende berekeningen --- en dat is
het praktische nut: welke kant van de identiteit van Green ook
gemakkelijker is, wordt de berekening, en de andere de controle.
Voor circulaties van veeltermvelden rond \omterm{def:b2:multint:exact}{gesloten} krommen is de
dubbele integraal bijna altijd de gemakkelijke kant.
\end{example}

\begin{remark}
Green--Riemann verklaart \cref{ex:b2:multint:angleform}: voor een
gesloten vorm ($Q_x = P_y$) verdwijnt de integraal rond de rand van
elk gebied dat in $U$ bevat is. De hoekvorm faalt alleen omdat het
gat in de oorsprong verhindert dat de schijf begrensd door de
eenheidscirkel binnen $U$ ligt --- lijnintegralen van gesloten
vormen sporen de gaten van het gebied op. (Doorgedreven wordt deze
waarneming de cohomologie van de Rham.)
\end{remark}

\section{Drievoudige integralen}

De theorie breidt zich zonder enig nieuw idee uit tot drie
veranderlijken: Fubini herleidt $\iiint$ tot drie integralen in één
veranderlijke (ofwel door \emph{in schijven te snijden}: $\iiint_K
f = \int\bigl(\iint_{K_z} f\bigr)\dd z$ over de horizontale
schijven $K_z$, ofwel door \emph{te stapelen}: eerst naar $z$
integreren langs verticale staafjes), en de formule voor de
verandering van veranderlijken geldt met de
$3 \times 3$-jacobimatrix.

\begin{example}[Cilinder- en bolcoördinaten]\label{ex:b2:multint:spherical}
\emph{Cilindrisch} $(x, y, z) = (\rho\cos\alpha, \rho\sin\alpha,
z)$: $\dd x\,\dd y\,\dd z = \rho\,\dd\rho\,\dd\alpha\,\dd z$.
\emph{Bolvormig} $(x, y, z) = (r\cos\theta\cos\varphi,\
r\sin\theta\cos\varphi,\ r\sin\varphi)$ ($\theta$ de lengtegraad,
$\varphi \in [-\frac\pi2, \frac\pi2]$ de breedtegraad): de $3
\times 3$-determinant langs de laatste rij ontwikkelen geeft
\[
\det J = r^2\cos\varphi ,
\qquad
\dd x\,\dd y\,\dd z
= r^2\cos\varphi\;\dd r\,\dd\theta\,\dd\varphi .
\]
\omterm{pb:b2:multint:1}{Volume van de bol} met straal $R$:
\[
V = \int_0^R\!\!\int_0^{2\pi}\!\!\int_{-\pi/2}^{\pi/2}
r^2\cos\varphi\;\dd\varphi\,\dd\theta\,\dd r
= \frac{R^3}{3}\cdot 2\pi \cdot 2
= \boxed{\frac43\pi R^3} ,
\]
waarmee eindelijk de formule wordt gekweten die in de
volumehoofdstukken van de eerdere boeken werd aangenomen.
\end{example}

\begin{example}[Het viervlak, tweemaal]
\label{ex:b2:multint:tetra}
Het volume van $T = \{x, y, z \geq 0,\ x + y + z \leq 1\}$, door
stapelen: voor vaste $(x, y)$ in de driehoek $x + y \leq 1$ loopt
$z$ over $\intcc0{1 - x - y}$, dus
\[
V = \int_0^1\!\!\int_0^{1-x}(1 - x - y)\,\dd y\,\dd x
= \int_0^1\frac{(1 - x)^2}{2}\,\dd x = \frac16 .
\]
Door in schijven te snijden: de doorsnede op hoogte $z$ is de
driehoek $\{x, y \geq 0,\ x + y \leq 1 - z\}$, met \omterm{def:b2:multint:domain}{oppervlakte}
$\frac{(1-z)^2}2$, en $V = \int_0^1\frac{(1-z)^2}2\,\dd z =
\frac16$ opnieuw --- de twee berekeningen zijn dezelfde integralen
in een andere volgorde, en meer beweert Fubini niet. De waarde
$\frac16 = \frac13\cdot\frac12\cdot1$ is de kegelformule
(\cref{ex:b2:multint:cone}) met driehoekige grondvlak, en de
$n$-dimensionale versie $1/n!$ wordt in de weekendopgave met
precies dit snijden bewezen.
\end{example}

\begin{example}[Zwaartepunt van een halve bol]\label{ex:b2:multint:centroid}
Voor de bovenste halve bol $H$ met straal $R$ ($z \geq 0$) is de
hoogte van het zwaartepunt $\bar z = \frac1{V}\iiint_H z$, met $V =
\frac23\pi R^3$. In bolcoördinaten ($z = r\sin\varphi$, $\varphi
\in \intcc0{\pi/2}$):
\[
\iiint_H z
= \int_0^R r^3\,\dd r\int_0^{2\pi}\dd\theta
\int_0^{\pi/2}\sin\varphi\cos\varphi\,\dd\varphi
= \frac{R^4}4\cdot2\pi\cdot\frac12 = \frac{\pi R^4}4 ,
\]
dus
\[
\bar z = \frac{\pi R^4/4}{2\pi R^3/3} = \frac{3R}8 :
\]
het evenwichtspunt van een massieve halve bol ligt op drie achtste
van de straal boven het platte vlak --- onder de halve hoogte
$R/2$, zoals het moet, omdat het lichaam nabij het grondvlak dikker
is. Elke berekening van een zwaartepunt heeft deze gedaante: één
momentintegraal, één volume, één verhouding, en een controle op
aannemelijkheid tegen de meetkunde.
\end{example}

\begin{example}[Volume door in schijven te snijden: de
kegel]\label{ex:b2:multint:cone}
Een kegel met grondvlak van \omterm{def:b2:multint:domain}{oppervlakte} $A$ en hoogte $h$ (top
boven, grondvlak op $z = 0$): de schijf op hoogte $z$ is het
grondvlak geschaald met de factor $(1 - z/h)$, met \omterm{def:b2:multint:domain}{oppervlakte} $A(1
- z/h)^2$. Bijgevolg is
\[
V = \int_0^h A\Bigl(1 - \frac zh\Bigr)^2\dd z = \frac{Ah}{3} :
\]
het één derde van de schoolformules, geldig voor \emph{elke} vorm
van het grondvlak --- in schijven snijden maakt er de integraal van
een kwadraat van.
\end{example}

\begin{example}[Drempels voor integreerbaarheid in het
vlak]\label{ex:b2:multint:threshold}
Voor welke $\alpha > 0$ convergeert $\iint_{D}\rho^{-\alpha}\,\dd
x\,\dd y$ op de geperforeerde eenheidsschijf $D$ (als limiet over
ringen $\varepsilon \leq \rho \leq 1$)? In poolcoördinaten is
\[
\int_0^{2\pi}\!\!\int_\varepsilon^1\rho^{-\alpha}\,
\rho\,\dd\rho\,\dd\alpha
= 2\pi\int_\varepsilon^1\rho^{1-\alpha}\,\dd\rho ,
\]
wat convergeert als $\varepsilon \to 0$ dan en slechts dan als $1 -
\alpha > -1$, dat wil zeggen $\alpha < 2$: in dimensie $2$ is de
kritieke exponent van de singulariteit de dimensie zelf, waarbij de
extra $\rho$ uit de jacobiaan de singulariteit met één macht
verzacht. (Evenzo $\alpha < 3$ voor een puntsingulariteit in de
ruimte, via $r^2$.) Deze radiale boekhouding is hoe de
integreerbaarheid in het kader van Lebesgue van bachelorjaar 3 in
één oogopslag wordt beslist --- en zij is de reden waarom $\iiint
1/r$ in \cref{exo:b2:multint:7} moeiteloos convergeerde.
\end{example}

\begin{remark}[Klassieke valkuilen]
(i) \emph{Oriëntatie}: een \omterm{def:b2:multint:lineint}{lijnintegraal} verandert van teken met de
richting van het doorlopen, en Green--Riemann vereist de rand tegen
de wijzers van de klok in (het gebied links); voor een gebied met
een gat wordt de binnenrand \emph{met} de wijzers van de klok mee
doorlopen. (ii) \emph{De jacobiaan komt binnen met een absolute
waarde}: een verandering van veranderlijken brengt nooit een
negatieve \omterm{def:b2:multint:domain}{oppervlakte} voort, en $\abs{\det}$ vergeten keert
gewoonlijk precies de tekens om wanneer de afbeelding de oriëntatie
omkeert. (iii) \emph{De poolfactor $\rho$}: $\dd x\,\dd y =
\rho\,\dd\rho\,\dd\alpha$, niet $\dd\rho\,\dd\alpha$ --- de meest
voorkomende fout van het hele hoofdstuk; dimensieanalyse vangt haar
op, want $\dd\rho\,\dd\alpha$ heeft de dimensie van een \omterm{def:b2:curves:length}{lengte}, niet
van een \omterm{def:b2:multint:domain}{oppervlakte}. (iv) \emph{Oneigenlijke dubbele integralen}:
de limieten over groeiende schijven en over groeiende vierkanten
stemmen hier overeen omdat de integranden positief zijn (inklemmen);
voor integranden die van teken wisselen kan de limiet van de
uitputting afhangen, en zonder \omterm{def:b2:series:def}{absolute convergentie} wordt niets
beweerd. (v) \emph{Gebieden tegenover integranden}: een
productintegrand op een gebied dat geen product is, ontbindt de
integraal \emph{niet} --- de ontbinding heeft beide nodig, zoals bij
het vierkant van \cref{ex:b2:multint:polar}.
\end{remark}

\begin{remark}[Vooruitblik binnen dit volume]
De integraal van Gauss die hier wordt berekend, zit stilzwijgend
overal in de kanshoofdstukken: de constante $\sqrt\pi$ binnen de
formule van Stirling (\cref{thm:b2:comparison:stirling}) is de
integraal van dit hoofdstuk, en via Stirling legt zij de asymptoten
$1/\sqrt{\pi n}$ vast van de terugkeerkansen van de toevalswandeling
in \cref{ch:b2:proba}. Ook de \omterm{pb:b2:multint:1}{integralen van Wallis} uit de
weekendopgave duiken daar op en drijven dezelfde schattingen van de
centrale binomiaalcoëfficiënten aan. In de andere richting maken de
oppervlakte- en volume-elementen van dit hoofdstuk de meetkunde van
\cref{ch:b2:surfaces} af, en de \omterm{pb:b2:multint:1}{formule van Green} berekent de
\omterm{def:b2:surfaces:area}{oppervlakten} van de \omterm{pb:b2:curves:1}{omhullenden} van \cref{ch:b2:curves} opnieuw (de
\omterm{pb:b2:curves:1}{astroïde}, in \cref{exo:b2:multint:5}). Eén hoofdstuk, drie
diensten: maat voor de meetkunde, constanten voor de kansrekening,
en de discipline van de verandering van veranderlijken voor beide.
\end{remark}

\section{Oefeningen}

\begin{exercise}[$\star$]\label{exo:b2:multint:1}
Bereken $\int_\gamma y^2\,\dd x + x\,\dd y$ langs: (a) het segment
van $(0,0)$ naar $(1,1)$; (b) de paraboolboog $y = x^2$ van $(0,0)$
naar $(1,1)$. Is de vorm \omterm{def:b2:multint:exact}{exact}?
\end{exercise}

\begin{exercise}[$\star$]\label{exo:b2:multint:2}
Toon aan dat $\omega = (2xy + y^3)\,\dd x + (x^2 + 3xy^2 + 1)\,\dd
y$ \omterm{def:b2:multint:exact}{gesloten} is op $\R^2$, bepaal een \omterm{def:b2:multint:exact}{potentiaal}, en bereken
$\int_\gamma\omega$ langs een willekeurige boog van $(0, 0)$ naar
$(1, 2)$.
\end{exercise}

\begin{exercise}[$\star$]\label{exo:b2:multint:3}
Bereken $\iint_D (x + y)\,\dd x\,\dd y$, waarbij $D$ het gebied is
begrensd door $y = x^2$ en $y = x$ ($0 \leq x \leq 1$), in beide
volgorden van integratie.
\end{exercise}

\begin{exercise}[$\star\star$]\label{exo:b2:multint:4}
Bereken met poolcoördinaten $\iint_D \frac{\dd x\,\dd y}{(1 + x^2 +
y^2)^2}$ over het hele vlak (als limiet over schijven), en
$\iint_{D'} xy\,\dd x\,\dd y$ over de kwartschijf $D' = \{x, y \geq
0,\ x^2 + y^2 \leq 1\}$.
\end{exercise}

\begin{exercise}[$\star\star$]\label{exo:b2:multint:5}
Bereken de \omterm{def:b2:multint:domain}{oppervlakte} omsloten door de \omterm{pb:b2:curves:1}{astroïde} $x = \cos^3 t$, $y
= \sin^3 t$, $t \in [0, 2\pi]$, met \cref{cor:b2:multint:area}.
\emph{(Lineariseer $\sin^2 t\cos^2 t$.)}
\end{exercise}

\begin{exercise}[$\star\star$]\label{exo:b2:multint:6}
Bereken het volume van het lichaam dat onder wordt begrensd door de
paraboloïde $z = x^2 + y^2$ en boven door het vlak $z = 1$, met
beide methoden: stapelen (integreer $1 - x^2 - y^2$ over de
eenheidsschijf, poolcoördinaten) en in schijven snijden (de
horizontale schijven zijn schijven met straal $\sqrt z$).
\end{exercise}

\begin{exercise}[$\star\star$]\label{exo:b2:multint:7}
(Zwaartekrachtige aantrekking van een bol --- de stelling van
Newton, bijzonder geval) Toon aan dat het volume van de bolschil $a
\leq r \leq b$ gelijk is aan $\frac43\pi(b^3 - a^3)$ en bereken
$\iiint_{B} \frac{\dd x\,\dd y\,\dd z}{r}$ over de bol $B$ met
straal $R$ ($r$ de afstand tot de oorsprong). \emph{(Bolcoördinaten.)}
\end{exercise}

\begin{exercise}[$\star\star\star$]\label{exo:b2:multint:8}
(Lemma van Poincaré, stervormig geval) Zij $U$ stervormig ten
opzichte van $0$ (dat wil zeggen $M \in U \Rightarrow [0, M]
\subseteq U$) en $\omega = P\dd x + Q\dd y$ een \omterm{def:b2:multint:exact}{gesloten}
$\mathcal{C}^1$-vorm op $U$. Definieer
\[
f(x, y) = \int_0^1 \bigl(x\,P(tx, ty) + y\,Q(tx, ty)\bigr)\dd t .
\]
Toon met differentiatie onder het integraalteken
(\cref{ch:b2:integration}) en $P_y = Q_x$ aan dat $f_x = P$ en $f_y
= Q$: elke gesloten vorm op een stervormige \omterm{def:b2:metric:topology}{open} verzameling is
\omterm{def:b2:multint:exact}{exact}.
\end{exercise}

\begin{exercise}[$\star\star\star$]\label{exo:b2:multint:9}
(De integraal van Dirichlet via dubbele integratie) Verantwoord en
buit uit
\[
\int_0^\infty\!\!\int_0^\infty e^{-xy}\sin x\;\dd y\,\dd x
\quad\text{tegenover}\quad
\int_0^\infty\!\!\int_0^\infty e^{-xy}\sin x\;\dd x\,\dd y
\]
op $[0, A] \times [0, \infty)$: toon aan dat $\int_0^A \frac{\sin
x}{x}\dd x = \frac\pi2 - \int_0^\infty e^{-Ay}\frac{y\sin A +
\cos A}{1 + y^2}\dd y$ en vind $\int_0^\infty \frac{\sin x}{x}\,\dd
x = \frac\pi2$ terug, in vergelijking met het bewijs via
\omterm{thm:b2:integration:continuity}{parameterintegralen} uit \cref{ch:b2:integration}.
\end{exercise}

\begin{exercise}[$\star\star\star$]\label{exo:b2:multint:10}
(Isoperimetrische ongelijkheid via Wirtinger) Zij $\gamma$ een
enkelvoudige \omterm{def:b2:multint:exact}{gesloten} $\mathcal{C}^1$-kromme met \omterm{def:b2:curves:length}{lengte} $2\pi$,
geparametriseerd naar \omterm{def:b2:curves:length}{booglengte} op $[0, 2\pi]$, die de \omterm{def:b2:multint:domain}{oppervlakte}
$A$ omsluit. Bewijs met \cref{cor:b2:multint:area}, Parseval en de
ongelijkheid van Wirtinger (oefeningen van \cref{ch:b2:fourier})
dat $A \leq \pi$, met gelijkheid voor de cirkel. \emph{(Normaliseer
$\int_0^{2\pi} x(s)\dd s = 0$; schrijf $2A = \oint x\,\dd y - y\,\dd
x$ en begrens $2A \leq \int (x^2 + y'^2)$ zorgvuldig via $2A =
\int_0^{2\pi}(xy' - yx')\dd s$ en $x^2 + y'^2 \geq 2xy'$.)}
\end{exercise}

\begin{exercise}[$\star\star$]\label{exo:b2:multint:11}
(Momenten van de bol) Bereken voor de bol $B$ met straal $R$ in
$\R^3$ de integraal $\iiint_B z^2\,\dd x\,\dd y\,\dd z$ in
bolcoördinaten, en leid met de symmetrie $\iiint_B (x^2 + y^2 +
z^2)\,\dd x\,\dd y\,\dd z$ af. Toets de laatste aan de berekening
met schillen $\int_0^R r^2\cdot4\pi r^2\,\dd r$.
\end{exercise}

\begin{exercise}[$\star\star$]\label{exo:b2:multint:12}
(Veterformule) Zij $K$ een veelhoek met hoekpunten $(x_1, y_1),
\dots, (x_m, y_m)$ tegen de wijzers van de klok in geordend
(indices modulo $m$). Leid uit \cref{cor:b2:multint:area} af dat
\[
\operatorname{Area}(K)
= \frac12\sum_{i=1}^m
\bigl(x_iy_{i+1} - x_{i+1}y_i\bigr) ,
\]
en ga de formule na op de driehoek $(0,0)$, $(1,0)$, $(0,1)$.
\end{exercise}

\section{Probleem: het volume van de bol in dimensie $n$}

\begin{problem}[{Weekendopgave --- $V_n =
\pi^{n/2}/\Gamma(\frac n2 + 1)$, en de vreemdheid van hoge
dimensies}]\label{pb:b2:multint:1}
De schijf heeft \omterm{def:b2:multint:domain}{oppervlakte} $\pi$, de bol volume $\frac43\pi$ ---
en dan? Deze opgave berekent het volume van de eenheidsbol van
$\R^n$ voor elke $n$, tweemaal (met een recursie door in schijven
te snijden, aangedreven door de \omterm{pb:b2:multint:1}{integralen van Wallis}, en daarna
met de $\Gamma$-functie en de integraal van Gauss uit
\cref{ex:b2:multint:polar}), en leest daarna de meetkunde af: de
volumes pieken in dimensie vijf en snellen naar nul, en bijna heel
een hoogdimensionale bol verbergt zich in een dunne schil nabij
haar rand.\index{volume van de bol}\index{integralen van Wallis}
Voor een \omterm{def:b2:metric:continuity}{continue} functie op een bol van $\R^n$ wordt de integraal
opgevat als de $n$-voudig herhaalde integraal (één coördinaat
tegelijk in schijven gesneden, als in het hoofdstuk voor $n \leq
3$); wij schrijven $B_n(R)$ voor de \omterm{def:b2:multint:exact}{gesloten} bol met straal $R$ en
middelpunt $0$, $v_n(R)$ voor haar volume, en $V_n = v_n(1)$, met
$V_0 = 1$ per afspraak.

\medskip
\noindent\textbf{Deel I --- De recursie door in schijven te
snijden.}
\begin{enumerate}
  \item Toon met de substitutie $x_i = Ru_i$ in elk van de $n$
        herhaalde integralen aan dat $v_n(R) = V_nR^n$.
  \item Toon, door $B_n(1)$ langs haar laatste coördinaat in
        schijven te snijden, aan dat
        \[
        V_n = V_{n-1}\int_{-1}^{1}(1 -
        t^2)^{\frac{n-1}2}\,\dd t .
        \]
  \item Identificeer met $t = \sin\theta$ de integraal als een
        integraal van Wallis: $\int_{-1}^1(1 -
        t^2)^{\frac{n-1}2}\dd t = 2W_n$, met $W_n =
        \int_0^{\pi/2}\cos^n\theta\,\dd\theta =
        \int_0^{\pi/2}\sin^n\theta\,\dd\theta$.
  \item Bewijs de twee identiteiten van Wallis (partieel
        integreren; dan $nW_nW_{n-1}$ telescoperen):
        \[
        W_n = \frac{n-1}nW_{n-2} \quad (n \geq 2),
        \qquad
        W_nW_{n-1} = \frac{\pi}{2n} \quad (n \geq 1).
        \]
\end{enumerate}

\medskip
\noindent\textbf{Deel II --- De recursie opgelost.}
\begin{enumerate}[resume]
  \item Combineer de vragen 2--4 tot de recursie met twee stappen
        \[
        V_n = \frac{2\pi}{n}\,V_{n-2}
        \qquad (n \geq 2).
        \]
  \item Leid de gesloten vormen af, voor $k \geq 0$:
        \[
        V_{2k} = \frac{\pi^k}{k!},
        \qquad
        V_{2k+1} = \frac{2^{k+1}\pi^k}{1\cdot3\cdot5\cdots
        (2k+1)} .
        \]
  \item Zet $V_1, \dots, V_7$ numeriek in een tabel. Bewijs met de
        verhouding $V_n/V_{n-2} = 2\pi/n$ en de waarden van $2W_5$
        en $2W_6$ dat de rij $(V_n)$ stijgt tot haar maximum $V_5 =
        \frac{8\pi^2}{15} \approx 5.26$ en daarna daalt.
  \item Toon aan dat $V_n \to 0$ sneller dan elke meetkundige rij,
        en dat $\sum_{n\geq1} V_n$ convergeert: alle eenheidsbollen
        samen hebben een eindig totaal volume.
  \item Bewijs de voortbrengende identiteit
        \[
        \sum_{k\geq0} V_{2k}\,x^{2k} = \eu^{\pi x^2}
        \qquad (x \in \R),
        \]
        en leid af dat $\sum_{k\geq0}V_{2k} = \eu^\pi \approx
        23.14$.
\end{enumerate}

\medskip
\noindent\textbf{Deel III --- Tweede weg: $\Gamma$ en de
integraal van Gauss.}
\begin{enumerate}[resume]
  \item Toon met Fubini aan (de integrand is een product) dat
        \[
        I_n = \int_{\R^n}\eu^{-\norm
        x^2}\dd x
        = \Bigl(\int_{-\infty}^{+\infty}
        \eu^{-t^2}\dd t\Bigr)^{\!n} = \pi^{n/2},
        \]
        de $n$-dimensionale integraal van Gauss, opgevat als een
        limiet over kubussen $\intcc{-R}{R}^n$.
  \item Herinner je $\Gamma(s) = \int_0^\infty t^{s-1}\eu^{-t}\dd
        t$ (\cref{def:b2:integration:gamma}). Bereken uit
        $\Gamma(s+1) = s\,\Gamma(s)$
        (\cref{thm:b2:integration:gammaprops}) en
        $\Gamma(\tfrac12) = \sqrt\pi$ (substitueer $t = u^2$ en
        roep de integraal van Gauss in):
        \[
        \Gamma(k + 1) = k!,
        \qquad
        \Gamma\Bigl(k + \frac32\Bigr) =
        \frac{1\cdot3\cdots(2k+1)}{2^{k+1}}\,\sqrt\pi .
        \]
  \item Bewijs met inductie via de recursie van vraag 5 de ene
        formule
        \[
        V_n = \frac{\pi^{n/2}}{\Gamma\bigl(\frac n2 +
        1\bigr)} \qquad (n \geq 1),
        \]
        en ga na dat zij beide gesloten vormen van vraag 6
        reproduceert.
  \item Toon aan dat $\int_0^\infty \eu^{-r^2}r^{n-1}\dd r =
        \tfrac12\Gamma\bigl(\tfrac n2\bigr)$ en leid de identiteit
        \[
        I_n = n\,V_n\int_0^\infty \eu^{-r^2}\,r^{n-1}\,\dd r
        \]
        af. Interpreteer haar: de gaussische massa van $\R^n$ wordt
        verzameld langs bolschillen waarvan de
        ``$(n-1)$-dimensionale \omterm{def:b2:multint:domain}{oppervlakte}'' op straal $r$ gelijk
        is aan $nV_nr^{n-1}$ --- beide leden zijn nu onafhankelijk
        bewezen, dus de interpretatie kost niets.
  \item Zet $s_{n-1} = nV_n$ (de \omterm{def:b2:multint:domain}{oppervlakte} van de eenheidssfeer
        $S^{n-1}$, in overeenstemming met $v_n(R) = \int_0^R
        s_{n-1}r^{n-1}\dd r$). Zet $s_0, \dots, s_3$ in een tabel
        en ga na dat $s_1 = 2\pi$, $s_2 = 4\pi$, $s_3 = 2\pi^2$.
\end{enumerate}

\medskip
\noindent\textbf{Deel IV --- Hoge dimensies zijn vreemd.}
\begin{enumerate}[resume]
  \item Toon met de formule van Stirling
        (\cref{thm:b2:comparison:stirling}) toegepast op $k!$ aan
        dat voor even $n = 2k$
        \[
        V_n \sim \frac{1}{\sqrt{\pi n}}
        \Bigl(\frac{2\pi\eu}{n}\Bigr)^{n/2}
        \qquad (n \to \infty, \ n \text{ even}),
        \]
        en leg uit waarom dezelfde grens van supermeetkundig verval
        zich via de recursie uitbreidt tot oneven $n$.
  \item De eenheidsbol ligt in de kubus $\intcc{-1}1^n$ met volume
        $2^n$. Bereken de vulverhouding $V_n/2^n$ voor $n = 2, 3,
        10$, en toon aan dat zij naar $0$ streeft: in hoge dimensie
        ligt vrijwel de hele kubus in haar hoeken.
  \item Toon aan dat het aandeel van $v_n(1)$ dat binnen afstand
        $\varepsilon$ van de randsfeer ligt, gelijk is aan $1 - (1
        - \varepsilon)^n \to 1$; welk aandeel van een bol in $100$
        dimensies ligt numeriek in de buitenschil met dikte
        $1\%$?
  \item Bewijs de asymptotiek van Wallis $W_n \sim
        \sqrt{\dfrac{\pi}{2n}}$ \emph{(de monotonie van $(W_n)$, de
        verhouding $W_n/W_{n-2} \to 1$, en $W_nW_{n-1} =
        \frac\pi{2n}$)}, en de ondergrens $W_n \geq
        \sqrt{\dfrac{\pi}{2(n+1)}}$ voor alle $n$.
  \item (Concentratie op een plak) Het aandeel van de eenheidsbol
        met eerste coördinaat voorbij $\delta$ is $\int_\delta^1(1
        - x^2)^{\frac{n-1}2}\dd x \,\big/\, (2W_n)$. Toon met $1 -
        u \leq \eu^{-u}$ en de staartgrens $\int_\delta^\infty
        \eu^{-a x^2}\dd x \leq \frac{\eu^{-a\delta^2}}{2a\delta}$
        aan dat dit aandeel hoogstens
        \[
        \frac{\eu^{-(n-1)\delta^2/2}}{(n-1)\,\delta}
        \Big/ \sqrt{\frac{2\pi}{n+1}}
        \]
        is, en besluit: voor $\delta = s/\sqrt{n-1}$ ligt op een
        aandeel $O(\eu^{-s^2/2}/s)$ na de hele bol in de plak
        $\abs{x_1} \leq s/\sqrt{n-1}$. Een bol van hoge dimensie is
        statistisch een dunne pannenkoek, in alle richtingen
        tegelijk.
  \item Zet de vragen 16--19 in elkaar tot één alinea: waar het
        volume van $B_n(1)$ zit (nabij de randsfeer, en toch binnen
        plakken van dikte $O(1/\sqrt n)$ van elk hypervlak door het
        middelpunt), en waarom deze twee uitspraken elkaar niet
        tegenspreken.
\end{enumerate}

\medskip
\noindent\textbf{Deel V --- Andere lichamen, en synthese.}
\begin{enumerate}[resume]
  \item (Simplex) Zij $\Delta_n = \{x \in \R^n : x_i \geq 0,\ \sum
        x_i \leq 1\}$. Toon met in schijven snijden en inductie aan
        dat $\operatorname{vol}(\Delta_n) = \frac1{n!}$.
  \item (Kruispolytoop) Leid af dat $C_n = \{x : \sum\abs{x_i} \leq
        1\}$ volume $\frac{2^n}{n!}$ heeft, en ga de insluiting
        $C_n \subseteq B_n(1) \subseteq \intcc{-1}1^n$ op het
        niveau van de volumes na: $\frac{2^n}{n!} \leq V_n \leq
        2^n$.
  \item Bereken $V_4$ op een derde manier: snijd $\R^4 = \R^2
        \times \R^2$ in schijven, integreer de \omterm{def:b2:multint:domain}{oppervlakte} van de
        $(z, w)$-schijf over de $(x, y)$-schijf in poolcoördinaten,
        en vind $V_4 = \frac{\pi^2}2$ terug.
  \item (Monte Carlo in moeilijkheden) Er wordt uniform een punt
        getrokken in de kubus $\intcc{-1}1^{20}$. Toon aan dat de
        kans dat het in de ingeschreven bol belandt gelijk is aan
        $V_{20}/2^{20} \approx 2.5\cdot10^{-8}$, zodat er ongeveer
        veertig miljoen trekkingen nodig zijn voordat de eerste
        treffer wordt verwacht: $V_n$ schatten met verwerpende
        steekproeven stort in hoge dimensie in (de vloek van de
        dimensionaliteit).
  \item Synthese. Twee onafhankelijke afleidingen ontmoetten elkaar
        in $V_n = \pi^{n/2}/\Gamma(\frac n2 + 1)$: som op welke
        stelling van dit hoofdstuk elk ervan gebruikte (Fubini, de
        verandering van veranderlijken, de integraal van Gauss in
        poolcoördinaten), en welke ingrediënten uit één
        veranderlijke (Wallis, $\Gamma$, Stirling). Waar doet het
        volume van bachelorjaar 3 deze berekening opnieuw met de
        theorie van Lebesgue, en wat voegt zij toe?
\end{enumerate}
\end{problem}
